\documentclass[10pt,a4paper]{amsart}
\usepackage{amsmath,amssymb,amsthm,graphicx}
% Retained preamble for the Stacks Project chapter snapshot in this directory.
% The chapter's own preamble.tex is a copy of the shared Stacks Project preamble; the block
% below is sliced MECHANICALLY, line for line, from it. Nothing was added, removed or
% redefined; the wrapper class, the xr-hyper/hyperref wiring and the externaldocument list
% are omitted because the assembler supplies the document class.
%
% lines taken from preamble.tex: 3, 5, 9, 11, 12, 156, 157, 158, ... (51 lines total)
%
% package coverage restored from the chapter's own preamble (2026-10-13): xy, xy-2cell

\usepackage[all]{xy}
\xyoption{2cell}
\UseAllTwocells

\newenvironment{reference}{\comment}{\endcomment}
\newenvironment{slogan}{\comment}{\endcomment}
\newenvironment{history}{\comment}{\endcomment}
\theoremstyle{plain}
\newtheorem{theorem}[subsection]{Theorem}
\newtheorem{proposition}[subsection]{Proposition}
\newtheorem{lemma}[subsection]{Lemma}
\theoremstyle{definition}
\newtheorem{definition}[subsection]{Definition}
\newtheorem{example}[subsection]{Example}
\newtheorem{exercise}[subsection]{Exercise}
\newtheorem{situation}[subsection]{Situation}
\theoremstyle{remark}
\newtheorem{remark}[subsection]{Remark}
\newtheorem{remarks}[subsection]{Remarks}
\numberwithin{equation}{subsection}
\def\lim{\mathop{\mathrm{lim}}\nolimits}
\def\colim{\mathop{\mathrm{colim}}\nolimits}
\def\Spec{\mathop{\mathrm{Spec}}}
\def\Hom{\mathop{\mathrm{Hom}}\nolimits}
\def\Ext{\mathop{\mathrm{Ext}}\nolimits}
\def\SheafHom{\mathop{\mathcal{H}\!\mathit{om}}\nolimits}
\def\SheafExt{\mathop{\mathcal{E}\!\mathit{xt}}\nolimits}
\def\Sch{\mathit{Sch}}
\def\Mor{\mathop{\mathrm{Mor}}\nolimits}
\def\Ob{\mathop{\mathrm{Ob}}\nolimits}
\def\Sh{\mathop{\mathit{Sh}}\nolimits}
\def\NL{\mathop{N\!L}\nolimits}
\def\CH{\mathop{\mathrm{CH}}\nolimits}
\def\proetale{{pro\text{-}\acute{e}tale}}
\def\etale{{\acute{e}tale}}
\def\QCoh{\mathit{QCoh}}
\def\Ker{\mathop{\mathrm{Ker}}}
\def\Im{\mathop{\mathrm{Im}}}
\def\Coker{\mathop{\mathrm{Coker}}}
\def\Coim{\mathop{\mathrm{Coim}}}
\DeclareMathSymbol{\boxtimes}{\mathbin}{AMSa}{"02}
\def\QCohstack{\mathcal{QC}\!\mathit{oh}}
\def\Cohstack{\mathcal{C}\!\mathit{oh}}
\def\Spacesstack{\mathcal{S}\!\mathit{paces}}
\def\Quotfunctor{\mathrm{Quot}}
\def\Hilbfunctor{\mathrm{Hilb}}
\def\Curvesstack{\mathcal{C}\!\mathit{urves}}
\def\Polarizedstack{\mathcal{P}\!\mathit{olarized}}
\def\Complexesstack{\mathcal{C}\!\mathit{omplexes}}
\def\Pic{\mathop{\mathrm{Pic}}\nolimits}
\def\Picardstack{\mathcal{P}\!\mathit{ic}}
\def\Picardfunctor{\mathrm{Pic}}
\def\Deformationcategory{\mathcal{D}\!\mathit{ef}}
\usepackage{cleveref}
\crefname{originaldefinition}{original definition}{original definitions}
\Crefname{originaldefinition}{Original definition}{Original definitions}
\makeatletter
\newlabel{definition-relative-representable-property}{{025V}{source}{}{sourcepointer-1}{}}
\newlabel{definition-relative-representable-property@cref}{{[originaldefinition][0][]025V}{[1][0][]source}{}{}{}}
\makeatother
\crefname{originallemma}{original lemma}{original lemmas}
\Crefname{originallemma}{Original lemma}{Original lemmas}
\makeatletter
\newlabel{lemma-morphism-schemes-gives-representable-transformation-property}{{02WJ}{source}{}{sourcepointer-2}{}}
\newlabel{lemma-morphism-schemes-gives-representable-transformation-property@cref}{{[originallemma][0][]02WJ}{[1][0][]source}{}{}{}}
\makeatother
\crefname{originallemma}{original lemma}{original lemmas}
\Crefname{originallemma}{Original lemma}{Original lemmas}
\makeatletter
\newlabel{lemma-representable-diagonal}{{025W}{source}{}{sourcepointer-3}{}}
\newlabel{lemma-representable-diagonal@cref}{{[originallemma][0][]025W}{[1][0][]source}{}{}{}}
\makeatother
\title{Fibre products of representable functors with a property make the diagonal representable with that property}
\author{The Stacks Project Authors}
\date{stacks-project:spaces.tex:local-snapshot-20261012; GNU FDL 1.2 or later}
\begin{document}
\maketitle
\setcounter{subsection}{1}
\noindent\textit{Editorial note.} All delivered mathematical text is the author's own, copied from the pinned Stacks Project chapter file spaces.tex; no statement or proof step was rewritten, reordered or removed.\par\medskip
\noindent\textit{Editorial note.} The proof contains one xymatrix display with two cartesian squares; it is part of the counted proof span and is preserved verbatim, with the chapter's own Xy-pic setup retained in the preamble.\par\medskip
\noindent\textit{Editorial note.} The three references of the counted spans are recorded as non-load-bearing navigation pointers; the argument that consumes them stays inside the counted proof span.\par\medskip
\noindent\textit{Editorial note.} Preamble: sources/stacks-spaces/retained-preamble.tex, sliced mechanically from the chapter's own preamble.tex.\par\medskip
\noindent\textit{Allowed standard prerequisites (editorial).}\begin{itemize}
\item Representability of the diagonal tested on pairs of maps. Lemma lemma-representable-diagonal (921-934), used at proof lines 973-974. Gives representability of Delta.
\item Morphisms of schemes correspond to representable transformations with the same property. Lemma lemma-morphism-schemes-gives-representable-transformation-property, used at proof lines 975-978. Reformulates condition (2) as a statement about transformations.
\item Stability of the property under base change. Definition definition-relative-representable-property, used at proof lines 990-994. Transfers property P to the base change of Delta.
\end{itemize}\medskip
\ifdefined\hypertarget\hypertarget{sourcepointer-1}{}\fi
\noindent\textit{Source reference.} Original definition 025V: focused-seeds/AGP/sources/stacks-snapshot-20261011/tags.txt, source lines 2828--2828. Navigation pointer into the Stacks Project; the standard fact it supplies is applied inside the counted proof span and the argument resting on it stays there.\par\medskip
\ifdefined\hypertarget\hypertarget{sourcepointer-2}{}\fi
\noindent\textit{Source reference.} Original lemma 02WJ: focused-seeds/AGP/sources/stacks-snapshot-20261011/tags.txt, source lines 3806--3806. Navigation pointer into the Stacks Project; the standard fact it supplies is applied inside the counted proof span and the argument resting on it stays there.\par\medskip
\ifdefined\hypertarget\hypertarget{sourcepointer-3}{}\fi
\noindent\textit{Source reference.} Original lemma 025W: focused-seeds/AGP/sources/stacks-snapshot-20261011/tags.txt, source lines 2829--2829. Navigation pointer into the Stacks Project; the standard fact it supplies is applied inside the counted proof span and the argument resting on it stays there.\par\medskip
\begin{lemma}
\label{lemma-transformation-diagonal-properties}
Let $S$ be a scheme contained in $\Sch_{fppf}$.
Let $F$ be a presheaf of sets on $(\Sch/S)_{fppf}$.
Let $\mathcal{P}$ be a property as in
Definition \ref{definition-relative-representable-property}.
If for every $U, V \in \Ob((\Sch/S)_{fppf})$ and $a \in F(U)$,
$b \in F(V)$ we have
\begin{enumerate}
\item $h_U \times_{a, F, b} h_V$ is representable, say by the scheme $W$, and
\item the morphism $W \to U \times_S V$ corresponding to
$h_U \times_{a, F, b} h_V \to h_U \times h_V$ has property $\mathcal{P}$,
\end{enumerate}
then $\Delta : F \to F \times F$ is representable and has
property $\mathcal{P}$.
\end{lemma}

\begin{proof}
Observe that $\Delta$ is representable by
Lemma \ref{lemma-representable-diagonal}.
We can formulate condition (2) as saying that
the transformation $h_U \times_{a, F, b} h_V \to h_{U \times_S V}$
has property $\mathcal{P}$, see Lemma
\ref{lemma-morphism-schemes-gives-representable-transformation-property}.
Consider $T \in \Ob((\Sch/S)_{fppf})$ and $(a, b) \in (F \times F)(T)$.
Observe that we have the commutative diagram
$$
\xymatrix{
F \times_{\Delta, F \times F, (a, b)} h_T \ar[d] \ar[r] &
h_T \ar[d]^{\Delta_{T/S}} \\
h_T \times_{a, F, b} h_T \ar[r] \ar[d] &
h_{T \times_S T} \ar[d]^{(a, b)} \\
F \ar[r]^\Delta & F \times F
}
$$
both of whose squares are cartesian. In this way we see that
the morphism $F \times_{F \times F} h_T \to h_T$ is the base
change of a morphism having property $\mathcal{P}$ by
$\Delta_{T/S}$. Since $\mathcal{P}$ is preserved under base
change this finishes the proof.
\end{proof}

\end{document}
